\documentclass[11pt,a4paper]{article}

\usepackage[T1]{fontenc}
\usepackage[utf8]{inputenc}
\usepackage{lmodern}
\usepackage[a4paper,margin=22mm]{geometry}
\usepackage{amsmath,amssymb,amsthm}
\usepackage{array}
\usepackage{graphicx}
\usepackage{systeme}
\usepackage{fancyhdr}

% Makrokomandos, naudotos originaliuose failuose.
\newcommand{\hm}[1]{\nobreak\hspace{0pt}#1\nobreak\hspace{0pt}}
\newcommand{\al}{\alpha}

% Uždavinių numeravimas.
\newcounter{uzdavinys}
\newenvironment{uzdavinys}
  {\refstepcounter{uzdavinys}\par\medskip\noindent\textbf{\theuzdavinys\ uždavinys.}\ }
  {\par\medskip}

% Jei brėžinio failo šalia nėra, dokumentas vis tiek susikompiliuos.
\makeatletter
\let\originalincludegraphics\includegraphics
\renewcommand{\includegraphics}[2][]{%
  \IfFileExists{#2}{\originalincludegraphics[#1]{#2}}{%
    \IfFileExists{#2.pdf}{\originalincludegraphics[#1]{#2}}{%
      \IfFileExists{#2.png}{\originalincludegraphics[#1]{#2}}{%
        \IfFileExists{#2.jpg}{\originalincludegraphics[#1]{#2}}{%
          \fbox{\parbox[c][28mm][c]{42mm}{\centering Brėžinys}}%
        }%
      }%
    }%
  }%
}
\makeatother

\setlength{\parindent}{0pt}
\setlength{\parskip}{0.25em}

\begin{document}


\begin{center}\LARGE\textbf{2016 m. IX--X klasių olimpiada}\end{center}

\section*{Uždavinių sąlygos}

\begin{center}
{\sc\large 2016 m.~9--10 klasės}
\end{center}
\vspace{0,3cm}

\setcounter{uzdavinys}{0}

\begin{uzdavinys}
Įrodykite, kad su visais teigiamais skaičiais $a$ ir $b$ teisinga nelygybė
\[
\frac{a+b}{1+a+b} < \frac{a}{1+a} + \frac{b}{1+b}.
\] 
\end{uzdavinys}

    
\vspace{0,25cm}



%\begin{uzdavinys}
\hspace*{-5.2mm}
\begin{minipage}%[b][3.2cm][c]
{0.79\textwidth}
\begin{uzdavinys}
Dešinėje pavaizduotą figūrą sudaro 5 vienodi kvadratai. Kiekvienas iš jų padalijamas į $n\times n$ vienetinių langelių. Reikia nudažyti penkis iš $5n^2$ langelių, vėl sudarančius figūrą ~\includegraphics[width=7mm]{./Breziniai/xbrezinys_2016}.
\begin{itemize}
\item[a)] Keliais būdais tai galima padaryti, jei $n=5$?
\item[b)] Kam lygu $n$, jei būdų tai padaryti iš viso yra 305?
\end{itemize}
\end{uzdavinys}
\end{minipage}
\quad\begin{minipage}%[t][1cm][t]
{0.8\textwidth}
\includegraphics[width=68.55pt]{./Breziniai/brezinys_2016}
\end{minipage}
%\end{uzdavinys}
\vspace{-1mm}



\begin{uzdavinys} 
Raskite mažiausią natūralųjį šešiaženklį skaičių, kurio visi skaitmenys yra skirtingi ir 
kuris dalijasi iš 2, 3, 4, 5 ir 6.
\end{uzdavinys}


\begin{uzdavinys} 
Duoti trys skirtingi natūralieji skaičiai. Bet kurių dviejų skaičių suma dalijasi 
iš likusio skaičiaus. 
\begin{itemize}
\item[a)] Raskite bent vieną tokį skaičių trejetą, kurio didžiausias skaičius būtų 9009.
\item[b)] Įrodykite, kad bet kurio tokio trejeto skaičių suma dalijasi iš 6.  
\end{itemize}
\end{uzdavinys}


\hspace*{-5.2mm}
\begin{minipage}%[b][3.2cm][c]
{0.78\textwidth}
\begin{uzdavinys}  
Stačiakampio formos popierinė juosta, kurios kraštinių ilgiai yra 2 cm ir 20 cm, perlenkiama taip, kad persidengiančios dalys sudaro trikampį, kaip parodyta paveikslėlyje. Kam lygus mažiausias galimas tokio trikampio plotas?
\end{uzdavinys} 
%\vspace{2cm}
\end{minipage}\quad\begin{minipage}%[t][1cm][t]
{0.8\textwidth}
\includegraphics[width=78.36pt]{./Breziniai/brezinys_2016_9_10}
\end{minipage}
%\vspace{0mm}
%\phantom{}

\vspace{0,7cm}

\clearpage

\section*{Sprendimai}

\begin{center}
{\sc\large 2016 m.~9--10 klasių uždavinių sprendimai}
\end{center}
\vspace{0,3cm}

\setcounter{uzdavinys}{0}

\begin{uzdavinys}
Įrodykite, kad su visais teigiamais skaičiais $a$ ir $b$ teisinga nelygybė
\[
\frac{a+b}{1+a+b} < \frac{a}{1+a} + \frac{b}{1+b}.
\] 
\end{uzdavinys}
       
\textit{Pirmasis sprendimas.} Tegu $a$ ir $b$ yra teigiami skaičiai. Kadangi $a>0$ ir $b>0$, tai
\[
\frac{a}{1+a+b} < \frac{a}{1+a} \quad \text{ir}\quad \frac{b}{1+a+b} < \frac{b}{1+b}.
\] 
Todėl
\[
\frac{a+b}{1+a+b} = \frac{a}{1+a+b} + \frac{b}{1+a+b} < \frac{a}{1+a} + \frac{b}{1+b}.
\]
 \vspace{0cm}  

\textit{Antrasis sprendimas.} Tegu $a$ ir $b$ yra teigiami skaičiai. Pertvarkykime duotąją nelygybę -- padauginkime ją iš teigiamo skaičiaus $(1+a+b)(1+a)(1+b)$ ir suprastinkime:
\begin{align*}
(a+b)(1+a)(1+b) &< a(1+b)(1+a+b) + b(1+a)(1+a+b),\\%[0,2cm]
(a+b)(1+a+b+ab) &< (1+a+b)(a(1+b) + b(1+a)),\\
a + a^2 + ab + a^2b + b + ab + b^2 + ab^2 &< (1+a+b)(a + b + 2ab),\\
a + a^2 + ab + a^2b + b + ab + b^2 + ab^2 &< a + b + 2ab + a^2 +ab +2a^2b +ab +b^2 + 2ab^2,\\
0 &< 2ab + a^2b + ab^2.
\end{align*}
Gautoji nelygybė teisinga, nes $a>0$, \,$b>0$. Todėl teisinga ir duotoji nelygybė.
\vspace{8mm} 

%\phantom{}

%\begin{uzdavinys}
\hspace*{-5.2mm}
\begin{minipage}%[b][3.2cm][c]
{0.79\textwidth}
\begin{uzdavinys}
Dešinėje pavaizduotą figūrą sudaro 5 vienodi kvadratai. Kiekvienas iš jų padalijamas į $n\times n$ vienetinių langelių. Reikia nudažyti penkis iš $5n^2$ langelių, vėl sudarančius figūrą ~\includegraphics[width=7mm]{./Breziniai/xbrezinys_2016}.
\begin{itemize}
\item[a)] Keliais būdais tai galima padaryti, jei $n=5$?
\item[b)] Kam lygu $n$, jei būdų tai padaryti iš viso yra 305?
\end{itemize}
\end{uzdavinys}
\end{minipage}
\quad\begin{minipage}%[t][1cm][t]
{0.8\textwidth}
\includegraphics[width=68.55pt]{./Breziniai/brezinys_2016}
\end{minipage}
\vspace{0cm}
       
\textit{Sprendimas.} 
Penkialangę figūrą, kurią reikia nudažyti, sudaro vidurinis langelis ir keturi langeliai, turintys su juo po bendrą kraštinę. Nustatykime, kiek langelių galima pasirinkti kaip tokios figūros vidurinį langelį.

Tinka visi langeliai, priklausantys viduriniam didžiajam kvadratui; jų yra $n^2$. Viršutiniame didžiajame kvadrate tinka visi langeliai, išskyrus viršutinę eilutę ir du kraštinius (iš kairės ir iš dešinės) stulpelius. Pašalinę šią eilutę ir stulpelius, vietoj $n\times n$ kvadrato gausime $(n-1)\times(n-2)$ stačiakampį, taigi $n^2-3n+2$ langelių. Po tiek pat tinkamų langelių dėl simetrijos yra ir likusiuose trijuose didžiuosiuose kvadratuose. Taigi tinkamų langelių, o tuo pačiu ir tinkamų penkialangės figūros nudažymo būdų iš viso~yra $$f(n)=n^2+4(n^2-3n+2)=5n^2-12n+8.$$
Konkretų tinkamų langelių skaičių $f(5)$ galima gauti ir neturint šios bendrosios $f(n)$ formulės, vietoj to tiesiogiai suskaičiavus pilkus langelius paveikslėlyje.

\begin{center}
\includegraphics[width=240pt]{./Breziniai/xxbrezinys_2016}
\end{center}

Dabar galime greitai atsakyti į uždavinio klausimus: a) $f(5)=73$; b) jei $f(n)\hm=305$, tai $5n^2-12n-297=0$ ir tada arba $n=9$, arba $n=-\frac{33}{5}$ (netinka, nes $n>0$).

\begin{flushright}
     Ats.: a) 73; b) 9.
    \end{flushright}
 \vspace{0,3cm}

  


\begin{uzdavinys} 
Raskite mažiausią natūralųjį šešiaženklį skaičių, kurio visi skaitmenys yra skirtingi ir 
kuris dalijasi iš 2, 3, 4, 5 ir 6.
\end{uzdavinys}

\textit{Sprendimas.} Natūralusis skaičius dalijasi iš 2 ir iš 5 tada ir tik tada, kai jis dalijasi iš~10, taigi kai baigiasi skaitmeniu 0. Tarp šešiaženklių skaičių $\overline{ABCDE0}$, kurių visi šeši skaitmenys skirtingi, mažiausi yra tie, kuriems $A=1$.
Tarp tokių skaičių $\overline{1BCDE0}$ -- tie, kuriems $B=2$. Tarp tokių skaičių $\overline{12CDE0}$ -- tie, kuriems $C=3$, o tarp tokių skaičių $\overline{123DE0}$ -- tie, kuriems $D=4$. 
Taigi mažiausi natūralieji šešiaženkliai skaičiai, kurių visi šeši skaitmenys skirtingi ir kurie dalijasi iš 2 bei iš 5, yra $\overline{1234E0}$, kur $E=5$, 6, 7, 8 arba 9. Reikia rasti mažiausią iš jų, kuris dalijasi iš 3, 4 ir 6 (jei tokio skaičiaus neatsirastų, turėtume nagrinėti didesnius skaičius). Skaičius $\overline{1234E0}=1234\cdot25\cdot4+10E$ dalijasi iš 3 ir iš 4, jei jo skaitmenų suma $10+E$ dalijasi iš 3, o skaičius $10E$ dalijasi iš~4. Tinka tik $E=8$. Gautasis skaičius $123\,480$ dalijasi iš 2, 5, 3, 4, taigi ir iš $6=2\cdot3$. Tai ir yra mažiausias skaičius, tenkinantis uždavinio sąlygą.

\begin{flushright}
     Ats.: $123\,480$.
\end{flushright}
\vspace{3mm}
 
 
 \begin{uzdavinys} 
Duoti trys skirtingi natūralieji skaičiai. Bet kurių dviejų skaičių suma dalijasi 
iš likusio skaičiaus. 
\begin{itemize}
\item[a)] Raskite bent vieną tokį skaičių trejetą, kurio didžiausias skaičius būtų 9009.
\item[b)] Įrodykite, kad bet kurio tokio trejeto skaičių suma dalijasi iš 6.  
\end{itemize}
\end{uzdavinys}

\textit{Sprendimas.} a) Uždavinio sąlygą tenkina skaičių trejetas 1, 2, 3. Jį radus, galima pastebėti, kad tinka ir visi trejetai $a$, $2a$, $3a$, kur $a$ yra bet koks natūralusis skaičius. Taip gauname tinkamą trejetą 3003, 6006, 9009, kurio didžiausias skaičius yra 9009.

b) Tarkime, kad natūralieji skaičiai $a$, $b$ ir $c$, kuriems $a<b<c$, tenkina uždavinio sąlygą. Kadangi 
$a+b$ dalijasi iš $c$, tai $a+b\geq c$. Kita vertus, $a+b<c+c=2c$, todėl $a+b=c$. Skaičius $a+c = 2a+b$ dalijasi iš $b$, todėl skaičius $2a$ dalijasi iš $b$. Kadangi 
$2a<2b$, tai $2a=b$. Taigi $c=a+b=3a$, o suma $a+b+c=6a$ dalijasi iš $6$. 

{\it Pastaba.} Čia iš tikrųjų įrodėme, kad be a)~dalyje pastebėtų trejetų $a$, $2a$, $3a$ kitų tinkamų skaičių trejetų nėra. Todėl gautasis a)~dalies atsakymas yra vienintelis galimas.
\begin{flushright}
     Ats.: a) 3003, 6006, 9009.
\end{flushright}
 \vspace{2mm}


\hspace*{-5.2mm}
\begin{minipage}%[b][3.2cm][c]
{0.78\textwidth}
\begin{uzdavinys}  
Stačiakampio formos popierinė juosta, kurios kraštinių ilgiai yra 2 cm ir 20 cm, perlenkiama taip, kad persidengiančios dalys sudaro trikampį, kaip parodyta paveikslėlyje. Kam lygus mažiausias galimas tokio trikampio plotas?
\end{uzdavinys} 
%\vspace{2cm}
\end{minipage}\quad\begin{minipage}%[t][1cm][t]
{0.8\textwidth}
\includegraphics[width=78.36pt]{./Breziniai/brezinys_2016_9_10}
\end{minipage}
\vspace{0mm}

\textit{Sprendimas.} Atkarpą, 
išilgai kurios perlenkta juosta, pažymėkime $AC$, o juostos~kraštų sankirtos tašką -- $B$. Trikampyje $ABC$ iš $A$ nuleiskime aukštinę $AH$ (žr.~pav.~kairėje). %\vfill

\begin{center}
\includegraphics[width=291.73pt]{./Breziniai/brezinys_2016_9_10_SPR}
%\vspace{0,2cm}
\end{center}

 Trikampio $ABC$ plotas $S$ lygus 
$\frac{1}{2}\cdot BC\cdot AH$. Atkarpos $AH$ ir $BC$ jungia popierinės juostos priešingus 
kraštus, ir jų ilgiai ne mažesni už juostos plotį 2 cm. Taigi 
\[
S=\frac{1}{2}\cdot BC\cdot AH\geq \frac{1}{2} \cdot 2 \cdot 2 =2 \text{ (cm}^2).
\]

Įmanoma taip perlenkti juostą, kad trikampio $ABC$ plotas būtų 2~cm$^2$. Iš tiesų, jei juostą perlenksime taip, kad $\angle ABC = 90^{\circ}$ (žr.~pav.~dešinėje), tai atkarpos $AB$ ir $BC$ bus lygios 
popierinės juostos pločiui, t.~y. $AB=BC=2$ cm. Tada $S=\frac{1}{2}\cdot BA\cdot BC\hm= \frac{1}{2} \cdot 2\cdot 2 =2 \text{ (cm}^2)$. Vadinasi, tai ir yra mažiausia $S$ reikšmė.

\begin{flushright}
     Ats.: 2 cm$^2$.
\end{flushright}

\vspace{0,7cm}


\end{document}
